\scanpage{101}
\mspage{698}

\diagram{\begin{tikzcd}[ampersand replacement=\&,column sep=small]
(234)\quad 1 \arrow[r] \&
\Pi_1(U_{0,3\,\overline{\mathbb Q}},-\bar{\jmath}) \arrow[r] \&
\Pi_1(U_{0,3\,\overline{\mathbb Q}},\mathfrak S_3;-\bar{\jmath}) \arrow[r] \&
\mathfrak S_3 \arrow[r] \& 1
\end{tikzcd}}

by the \deleted{\unsure{invariant}} commutator subgroup
of the group $\Pi_1(U_{0,3\,\overline{\mathbb Q}},-\bar{\jmath})$.
This \deleted{is} extension structure (233) of $\mathfrak S_3$
by $M^\wedge$ \uline{is not} \uline{trivial}, unless I am mistaken,
although the induced extension of
$\{1,\dot\rho,\dot\rho^2\}\simeq\mathbb Z/3\mathbb Z$ by $M^\wedge$,
or of $\{1,\dot\sigma\}$ by $\widehat M$, is.
\margin{yes, it is trivial, cf. below...}
A fortiori (232) is not trivial (although it has canonical
splittings \unsure{below} $\widehat{\mathbb Z}$, above
$\{1,\dot\rho,\dot\rho^2\}$ and above $\{1,\dot\sigma\}$).
Thus (subject to checking these points \unsure{presently}...)
the group $G_{\rho,\sigma}/\mu_G$ cannot be identified
\unsure{with a} subgroup of the \inserted{\unsure{affine|linear}} group
\[
(234)\qquad
\boxed{\operatorname{Aff}_{\widehat{\mathbb Z}}(\widehat M)
\simeq\operatorname{Aut}_{\widehat{\mathbb Z}}(\widehat M)\mathbin{.}\widehat M}
\]
corresponding to the subgroup $\widehat{\mathbb Z}^{*}\times\mathfrak S_3$
of $\operatorname{Aut}_{\widehat{\mathbb Z}}(\widehat M)$.

\deleted{\unsure{This|And}} Let us consider the \inserted{sub-}group
\[
(235)\qquad
\operatorname{Aff}\simeq
(\widehat{\mathbb Z}^{*}\times\mathfrak S_3)\mathbin{.}\widehat M
\]
We have a diagram
\diagram{\begin{tikzcd}[ampersand replacement=\&,column sep=small]
(237)\quad 1 \arrow[r,dash] \& M^\wedge \arrow[r,dash] \&
\operatorname{Aff} \arrow[r] \&
\widehat{\mathbb Z}^{*}\times\mathfrak S_3 \arrow[r,dash] \& 1 \\
\&\& \mathcal M_{0,3} \arrow[ur] \&\&
\end{tikzcd}}
with exact rows. The giving of the structure (230) is equivalent,
up to isomorphism (nonunique), to the

\margin{100}

\scanpage{102}

giving of an $\widehat M$-conjugacy class of liftings
of $\mathcal M_{0,3}\longrightarrow\mathfrak S_3\times\widehat{\mathbb Z}^{*}$
to a
\diagram{\begin{tikzcd}[ampersand replacement=\&]
(238)\quad\mathcal M_{0,3} \arrow[r] \& \operatorname{Aff}.
\end{tikzcd}}
\margin{\uline{N.B.} This homomorphism induces on $\Pi_{0,3}^\wedge$ the canonical homomorphism
\diagram{\begin{tikzcd}[ampersand replacement=\&,column sep=small]
\Pi_{0,3}^\wedge \arrow[r] \& (\Pi_{0,3}^\wedge)_{\mathrm{ab}}
\arrow[r,"\sim"] \& \widehat M
\end{tikzcd}}
an \uuline{essential} property of this homomorphism.}

The choice of such a lifting depends on the choice of a point of
$U_\infty(\overline{\mathbb Q})$ above the point $-\bar{\jmath}$
of $U_{0,3}(\overline{\mathbb Q})$ --- the $\mathrm{chg}^{t}$ of this
origin point is made by translation by an $x\in\widehat M$,
and the lifting (238) is then changed by \unsure{inner automorphism}
by this element $x$.

Thus the Fermat tower provides such a conjugacy class
of liftings\inserted{(238)}, knowledge of which is in principle equivalent
to that of the Fermat tower up to isomorphism. We would like to put,
if possible, these liftings in correspondence with
\unsure{the|the} homomorphisms (217), and more particularly
\unsure{one of them} (218), which provides a
\inserted{canonical} homomorphism from $\mathcal M_{0,3}$, not into
\deleted{\illegible{1}} $\operatorname{Aff}$, but into
\deleted{$\unsure{G_{\rho,\sigma}/\mu_G}$}
the analogous group $G_{\rho,\sigma}/\mu_G$ of (232), which lifts
the same homomorphism
$\mathcal M_{0,3}\longrightarrow
\deleted{\unsure{\mathbb O|\Phi}}\inserted{\widehat{\mathbb Z}}^{*}\times\mathfrak S_3$.

\margin{101}

\scanpage{103}
\mspage{699}

\diagram{\begin{tikzcd}[ampersand replacement=\&,column sep=small,row sep=small]
(239)\quad 1 \arrow[r] \& M^\wedge \arrow[r] \&
G_{\rho,\sigma}/\mu_G \arrow[r] \arrow[d,phantom,"\simeq"{rotate=90}] \&
\widehat{\mathbb Z}^{*}\times\mathfrak S_3 \arrow[r] \& 1. \\
\&\& \widehat{\mathbb Z}^{*}\mathbin{.}\mathfrak S^{+\wedge}/\mu_{\mathfrak S} \&\& \\
\&\&\& \mathcal M_{0,3}=\Pi_1(U_{0,3\,\mathbb Q},\mathfrak S_3;-\bar{\jmath})
\arrow[uu] \arrow[ul,dashed] \&
\end{tikzcd}}

On $\mathfrak S_{0,3}^{+\wedge},\simeq
\Pi_1(U_{0,3\,\overline{\mathbb Q}},\mathfrak S_3,-\bar{\jmath})
\subset\mathcal M_{0,3}$,
this homomorphism induces the canonical homomorphism
\deleted{of passage to the \unsure{quotient} $\unsure{\mathfrak S|G}$}
\inserted{from the} quotient of $\mathfrak S_{0,3}^{+\wedge}$
obtained ``by making $\Pi_{0,3}^\wedge$ abelian'',
to $G_{\rho,\sigma}^{\unsure{\natural|\sharp}}/\mu_G
\simeq\mathfrak S^{+\wedge}/\mu_{\mathfrak S}$, an extension of
$\mathfrak S_3$\inserted{$\subset\widehat{\mathbb Z}^{*}\times\mathfrak S_3$}
by $M^\wedge$.

Let us take up again the homomorphism $\mathcal M_{0,3}\longrightarrow\operatorname{Aff}$ (238),
it fits into an exact sequence
\diagram{\begin{tikzcd}[ampersand replacement=\&,column sep=small,row sep=large]
(240)\quad 1 \arrow[r] \& \Pi_{0,3}^\wedge \arrow[r] \arrow[d] \&
\mathcal M_{0,3} \arrow[r,dash] \arrow[d] \&
\overset{\inserted{(\simeq\deleted{\unsure{\mathfrak S_3}}\Gamma_{\mathbb Q}\times\mathfrak S_3}}
{\mathcal N_{\deleted{\unsure{\mathbb Q}}\unsure{0,3}}}
\arrow[r,dash] \arrow[d] \& 1 \\
1 \arrow[r] \& \widehat M \arrow[r] \& \operatorname{Aff} \arrow[r,dash] \&
\deleted{\unsure{\Gamma_{\mathbb Q}}}\inserted{\widehat{\mathbb Z}^{*}}
\times\mathfrak S_3 \arrow[r,dash] \& 1,
\end{tikzcd}}
where $\Gamma_{\mathbb Q}\times\mathfrak S_3\longrightarrow
\widehat{\mathbb Z}^{*}\times\mathfrak S_3$ is the homomorphism
\inserted{$\chi\times\mathrm{id}$} of (237), and where
$\Pi_{0,3}^\wedge\longrightarrow\widehat M$ is the canonical homomorphism
\diagram{\begin{tikzcd}[ampersand replacement=\&,column sep=small]
(241)\quad\Pi_{0,3}^\wedge \arrow[r] \&
(\Pi_{0,3}^\wedge)_{\mathrm{ab}} \arrow[r,phantom,"\simeq"] \&
\bigl((\Pi_{0,3})_{\mathrm{ab}}\bigr)^\wedge
\arrow[r,"\sim"] \& \widehat M.
\end{tikzcd}}
If we denote by \deleted{$\unsure{\mathcal M_{0,2}|\mathcal M_{0,3}}$}
$\mathcal M_1$ the quotient of $\mathcal M_{0,3}$ by
$[\Pi_{0,3}^\wedge,\Pi_{0,3}^\wedge]$, we find \unsure{therefore}

\margin{102}

\scanpage{104}

a \deleted{canonical} homomorphism \inserted{of sequences} that are exact
\diagram{\begin{tikzcd}[ampersand replacement=\&,column sep=small,row sep=large]
1 \arrow[r,dash] \& M^\wedge \arrow[r,dash] \arrow[d,Rightarrow] \&
\mathcal M_1 \arrow[r] \arrow[d] \&
\mathcal N_{\deleted{\unsure{\mathbb Q}}0,3} \arrow[r] \arrow[d] \& 1 \\
(242)\quad 1 \arrow[r,dash] \& M^\wedge \arrow[r] \&
\operatorname{Aff} \arrow[r,dash] \&
\widehat{\mathbb Z}^{*}\times\mathfrak S_3 \arrow[r,dash] \& 1,
\end{tikzcd}}
(depending on the choice of a $\mathrm p^{t}$ of $U_\infty(\overline{\mathbb Q})$
above $-\bar{\jmath}$)
which shows that the extension \inserted{$\mathcal M_1$} of
$\Gamma_{\mathbb Q}$ by $\widehat M$
\deleted{is identified}\inserted{is therefore}
\unsure{with|:} the inverse image of the trivial extension
$\operatorname{Aff}$ of $\widehat{\mathbb Z}^{*}\times\mathfrak S_3$
by $M^\wedge$, hence is trivial.
More precisely still, let us consider the quotient \inserted{$\mathcal M_2$}
of $\mathcal M_1$ by the kernel of $\mathcal M_1\longrightarrow\operatorname{Aff}$,
we find a diagram
\diagram{\begin{tikzcd}[ampersand replacement=\&,column sep=small,row sep=large]
1 \arrow[r,dash] \& M^\wedge \arrow[r] \arrow[d,dash,"\sim"{rotate=90}] \&
\mathcal M_2 \arrow[r] \arrow[d,"\sim"{rotate=90}] \&
N_2 \arrow[r] \arrow[d,"\sim"{rotate=90}] \& 1 \\
(242\ \mathrm{bis})\quad 1 \arrow[r,dash] \& M^\wedge \arrow[r] \&
\operatorname{Aff} \arrow[r,dash] \&
\widehat{\mathbb Z}^{*}\times\mathfrak S_3 \arrow[r,dash] \& 1
\end{tikzcd}}
(where $N_2\overset{\mathrm{def}}{\simeq}\mathcal M_2/\operatorname{Im}M^\wedge$).
We thus see that the extension of $N_2$ by $M^\wedge$ is trivial
--- we have an $M^\wedge$-conjugacy class of splittings of this
extension --- the splittings corresponding to the $\mathrm p^{\mathrm{ts}}$ of
\deleted{$\unsure{U_{0,2}|U_{0,3}}$}
$U_\infty(\overline{\mathbb Q})$ above the $\mathrm p^t$
$-\bar{\jmath}$ of $U_{0,3}(\overline{\mathbb Q})$

\margin{103}

\scanpage{105}
\mspage{700}

We would really like to show that this homomorphism
$\mathcal M_{0,3}\longrightarrow\operatorname{Aff}$ is isomorphic
to the homomorphism $\mathcal M_{0,3}\longrightarrow G_{\rho,\sigma}/\mu_G$
of (239), by means of a suitable isomorphism of extensions between
$\operatorname{Aff}$ and $G_{\rho,\sigma}/\mu_G$.
This would imply that this latter extension is split.
The isomorphisms of this extension onto $\operatorname{Aff}$
correspond \unsure{exactly} to the splittings of this extension.

$12^\circ)$ \uline{Splittings of the extension}
$\underline G_{\rho,\sigma}/\mu_G$ of
$\mathbb O_{12}^{*}\times\mathfrak S_3$
by $\underline M$.

Since we already have a lifting
\diagram{\begin{tikzcd}[ampersand replacement=\&]
\mathbb O_{12}^{*} \arrow[r,hook] \& \underline G_{\rho,\sigma}
\end{tikzcd}}
$(\unsure{150|180},151)$, I am going to look at the splittings
of the extension (231) which extend this one --- or what
\uline{amounts to the same thing}, the splittings
of the induced extension on $\mathfrak S_3$:
\margin{not quite --- one \unsure{would have to|has to} restrict oneself to the splittings
\diagram{\begin{tikzcd}[ampersand replacement=\&,column sep=small]
(\mathfrak S_3)_{S_0} \arrow[r] \& \underline{\mathfrak S}^{+}/\mu_{\mathfrak S}
\end{tikzcd}}
which \uline{commute with}
$\unsure{\underline G_{\rho,\sigma}(0)}\simeq
\mathbb O_{12}^{*}$}
\diagram{\begin{tikzcd}[ampersand replacement=\&,column sep=small]
(243)\quad 1 \arrow[r,dash] \& \underline M \arrow[r] \&
\underline{\mathfrak S}^{+}/\mu_{\mathfrak S} \arrow[r] \&
\underbrace{(\mathfrak S_3)_{S_0}}_{(H/\mu)_{S_0}} \arrow[r] \& 1
\end{tikzcd}}
\margin{We look at the \deleted{\unsure{compatible}}
splittings over an arbitrary base $S$.}
The splittings correspond \deleted{\unsure{is}}
bijectively to the pairs

\margin{104}

\scanpage{106}

\[
\left\{\begin{aligned}
 &(\rho',\sigma')=(\zeta\rho,\eta\rho)
   \qquad\text{with }\zeta,\eta\in\Gamma\underline M
   \;\deleted{\unsure{\widehat M}}\\
 &\text{satisfying}\\
 &\rho'{}^3=\sigma'{}^2=1,\qquad
   \sigma'(\rho')=\rho'{}^{-1}.
\end{aligned}\right.
\tag{244}
\]

We are going to interpret these conditions in terms of $\zeta,\eta$.
The condition $\rho'{}^3=1$ is written
(since $\unsure{\dot\rho|\rho}^3=1$ in $(\mathfrak S_3)_{S_0}$)
\[
\zeta\rho(\zeta)\rho^2(\zeta)=1
\qquad\text{i.e. additively}\qquad
(1+\rho_{\underline M}+\rho_{\underline M}^2)(\zeta)=0
\tag{245}
\]
this condition is \unsure{always|always} satisfied, for we have
\[
1+\rho_{\underline M}+\rho_{\underline M}^2=0,
\tag{\unsure{245}}
\]

The condition $\sigma'{}^2=1$ is written
\[
\eta\sigma(\eta)=1
\qquad\text{i.e. additively}\qquad
(1+\sigma_{\unsure{M|\underline M}})(\eta)=0
\tag{246}
\]
and it means that we have
\[
\eta=b\delta\qquad
\begin{gathered}
\text{with }\unsure{b|\beta}\in\Gamma\mathbb G_a.\\
\delta=\lambda_1-\lambda_0
\end{gathered}
\tag{247}
\]

We shall set
\[
\zeta=a_0\lambda_0+a_1\lambda_1
\qquad a_0,a_1\in\Gamma\mathbb G_a,
\tag{248}
\]
and we are going to express in terms of $a_0,a_1,b$ the condition
$\sigma'(\rho')=\rho'{}^{-1}$.

We have
\[
\sigma'(\rho')=\operatorname{int}(\eta)\sigma(\zeta\rho)
=\operatorname{int}(\eta)
\left[\sigma(\zeta)\underbrace{\sigma(\rho)}_{\rho^{-1}\lambda_1}\right]
\]
the condition
$\sigma'(\rho')=\underbrace{\rho'{}^{-1}}_{\rho^{-1}\zeta^{-1}}$
is written
\[
\sigma(\zeta)\rho^{-1}\lambda_1
\overset{\deleted{\unsure{=}}}{=}
\eta^{-1}\rho^{-1}\zeta^{-1}\eta
\]
i.e., multiplying on the \unsure{left|left} by $\rho$
\[
\rho(\sigma(\zeta))\mathbin{.}\lambda_1
=\rho(\eta)^{-1}\zeta^{-1}\eta
\]
or again, additively:

\margin{105}

\scanpage{107}
\mspage{701}

\[
\underbrace{\rho_M\sigma_M}_{\varepsilon_{1M}}(\zeta)+\lambda_1
=-\rho_M(\eta)-\zeta+\eta
\qquad\text{or again}
\]
\[
\underbrace{(1+(\varepsilon_1)_M)(\zeta)}_{(2a_1-a_0)\lambda_1}
+\lambda_1
=\underbrace{(1-\rho_M)(\eta)}_{3b\lambda_1},
\tag{249}
\]
or again, by an immediate calculation
\[
2a_1-a_0+1=3b
\qquad\text{or}\qquad
2a_1-a_0-3b\unsure{=|\ne}-1.
\tag{250}
\]

This shows that for $\zeta\inserted{\in\Gamma\underline M}$ given
\inserted{(i.e. for $\rho'$ given)}, for there to exist an $\eta$
(i.e. $\unsure{\sigma'|b}$) satisfying \unsure{with} $\zeta$ these
conditions (244) --- i.e. \deleted{\unsure{of|and}} such that
$(\rho'=\zeta\rho,\sigma'=\eta\sigma)$ \unsure{defines} a splitting of
(243) --- it is necessary and sufficient that we have
\[
2a_1-a_0+1\equiv0\qquad(3)
\tag{251}
\]
and then $\eta$ is uniquely determined by (250), at least if $3$ is
a regular section of $\underline{\mathcal O}_S$ on the base $S$
under consideration (which is the case if $S=\operatorname{Spec}\widehat{\mathbb Z}$).

Thus the splittings of (243) form a scheme
\inserted{$\underline{\mathrm{Sc}}$} over $S_0$, which is
\deleted{canonical}\inserted{given up to} \unsure{isomorphism} by
\[
\begin{aligned}
\underline{\mathrm{Sc}}
&\simeq
\{(a_0,a_1,b)\in\Gamma\mathbb E^3
  \mid 2a_1-a_0-3b=-1\}\\
&\simeq
\{(\zeta,\eta)\in\Gamma(M\times\underline{\mathbb Z\delta})
  \mid(1+\varepsilon_{M})(\zeta)
       -(1-\rho_M)(\eta)=-\lambda_1\}
\end{aligned}
\tag{252}
\]
\deleted{which is}

This scheme is naturally a

\margin{106}